\section{Targeted Procedure Intervention}
\label{sec:intervention}
Each frozen endpoint is replayed under three conditions while its evidence package, model configuration, prompt harness, evaluator, and period remain fixed. The intervention is intentionally simple: compare the original agent with the same agent receiving either a target-blind helper or the one reusable operation predicted by the registered failure mechanism.

\paragraph{Design requirements.}
The base comparison must satisfy three requirements. First, the targeted arm must improve over the original no-skill agent; otherwise there is no repaired failure. Second, it must also improve over a helper with the same callable surface and matched implementation complexity; otherwise the effect may reflect generic assistance rather than the target operation. Third, the endpoint outcome and transfer split must be frozen before the corresponding model outcomes are opened, so neither the mechanism label nor the success criterion can move to fit the result. We then add two stronger post-review requirements without changing the registered family outcome: G$_0$ must first pass a prospective neutrality test against fresh N, and R must expose the same operation output and candidate evidence as T while changing only the integration surface.

The targeted condition injects the registered mechanism-specific procedure. The generic condition injects a frozen target-blind helper under the same callable wrapper. The no-skill condition adds no experimental helper. The scientific unit is the frozen endpoint, not a repeat. We evaluate the two contrasts from Definition~1 jointly: targeted minus no skill asks whether the operation repairs the original agent, while targeted minus generic asks whether that repair exceeds generic helper/code assistance. Neither contrast is sufficient alone.

\subsection{Mechanism contracts}
The cutoff procedure filters candidate evidence by the target cutoff. Its generic control validates package structure without using timestamps. The release-alignment procedure normalizes the reporting period represented by a source release. Its generic control inspects document structure without performing date-to-period mapping. The grounding procedure enriches candidate documentary spans using reusable attribution and outlook cues. Its control can expose document structure but has no contextual event selector.

No targeted procedure contains scenario-specific values, dates, labels, or expected answers. Procedure source hashes are frozen before execution. The targeted and generic source implementations are matched exactly in lexical-token count and AST-node count for each family. This generic control is deliberately \emph{target-blind}, not assumed behaviorally neutral: it controls the callable wrapper and implementation-complexity surface while withholding the temporal operation under test. The no-skill arm is therefore essential for detecting harm introduced by the particular generic helper. G$_0$ addresses the remaining concern directly: it executes one helper call but is information-empty, reading neither package nor context and returning only an empty object. Its neutrality is tested empirically rather than assumed.

\subsection{Resource parity and identification layers}
The experiment measures controlled procedure exposure, not autonomous skill selection. Under the frozen harness, T, G, and G$_0$ each execute exactly one assigned helper operation before generation; N executes no experimental helper. R instead executes the same targeted temporal operation during context construction and exposes no callable helper. Thus T--G matches callable surface and source complexity, G$_0$ tests whether extra helper computation itself perturbs the agent, and T--R holds operation content fixed while changing only where that content enters the prompt. We make no latency or provider-cost superiority claim.

The completed identification ladder has four distinct contrasts. T--N asks whether the operation repairs the original agent. T--G asks whether it beats the executed target-blind helper. T--G$_0$ asks whether the operation retains value beyond behaviorally neutral extra computation, after G$_0$ has independently passed its N-equivalence gate. T--R asks whether callable skill state adds value beyond an operation-matched retrieval/context implementation. For R, the original candidate evidence is preserved and the exact T operation output is added as \texttt{retrieval\_operation\_output}; a zero-call audit verifies content equality and that no candidate evidence is removed or reordered. This is intentionally a strongest same-information baseline, not a weaker temporal-RAG surrogate.

\subsection{Frozen family outcomes and scorer timeline}
Before any source-native model outcome is opened, the failure-family contracts define what constitutes success. A cutoff success requires a nonempty set of cited evidence and a selected latest reference that are all available by the target cutoff. A release-alignment success requires the correct reporting-period identity together with the aligned headline value. A grounding success requires the requested evidence slots to be filled by frozen first-party sentence IDs in the correct slot order.

An auxiliary full-task conjunctive scorer is preserved only as sensitivity analysis because it adds orthogonal output fields and exact surface-form requirements beyond the registered mechanism outcome. The family definitions frozen before execution remain primary; Appendix~A records the scorer timeline and representation audit.

The base decision rule follows Definition~1: T$>$G with T$\leq$N is no repair, while T$>$N without T$>$G cannot isolate the targeted operation from generic assistance. The stronger controls are interpreted asymmetrically. G$_0$ can strengthen operation specificity only after its fresh G$_0$--N 90\% interval lies wholly inside the predeclared $\pm10$-point margin. For container value, a new callable-skill claim requires a prospectively resolved T--R residual of at least 10 points with the frozen endpoint-level criterion. If the 90\% T--R interval instead rules out a $>10$-point advantage, we retain only the operation-level claim and explicitly concede retrieval/context parity.

\subsection{Held-out transfer and evidence layers}
The initial replay freezes twelve pre-April endpoints as C3-R and eleven April--May endpoints as held-out C4-R; same-domain and compatible cross-domain labels are fixed before outcomes, and procedure source is unchanged across the split. A separate EIA validation freezes four earlier and eight held-out energy endpoints while reusing T1/G1 byte-identically. DeepSeek is the registered primary track, Kimi a provenance-clean replication selected by a pre-outcome support rule, and the mini-model layer an exploratory capacity stress. Full execution chronology, the failed local-secondary gate, and row-level provenance bindings are reported in Appendix~A.
